%! Vector-Valued Functions — Unit 4, Lesson 4.9 — Group Activity (Answer Key)
\documentclass[10pt]{article}
\usepackage{precalculus-article}
\usepackage{precalculus-key}   % pulls in -boxes; do NOT also load -boxes
\newcommand{\TallMath}[1]{$\displaystyle #1\rule[-1.4em]{0pt}{3.2em}$}

\begin{document}

\pageheader{Unit 4, Lesson 4.9}{Group Activity --- Answer Key}
\namepartnerperiod

\vspace{0.1in}

\begin{scenariobox}[Tracking a Particle]{sky}
\small
In each tier below you will work with vector-valued functions to describe the
position and motion of a particle in the plane.  Use vector notation
$\vec{p}(t)=\langle x(t),\,y(t)\rangle$ and the magnitude formula
$\|\vec{p}\|=\sqrt{[x(t)]^2+[y(t)]^2}$.
\end{scenariobox}

\vspace{0.1in}

%----------------------------------------------------------------------
% Tier R
%----------------------------------------------------------------------
\begin{tcolorbox}[colback=white, colframe=black!40,
    title=\textbf{Tier R --- Remediate --- Key},
    fonttitle=\bfseries, arc=2mm, left=3mm, right=3mm, top=2mm, bottom=2mm]
\small

$\vec{p}(t)=\langle 3t,\,2t\rangle$.

\vspace{0.06in}
\textbf{(a)}

\renewcommand{\arraystretch}{1.6}
\hspace{1em}
\begin{tabular}{|c|c|c|c|}
\hline
$t$ & $x(t)=3t$ & $y(t)=2t$ & $\vec{p}(t)$ \\
\hline
$0$ & \ans{$0$} & \ans{$0$} & \ans{$\langle 0,0\rangle$} \\
\hline
$1$ & \ans{$3$} & \ans{$2$} & \ans{$\langle 3,2\rangle$} \\
\hline
$2$ & \ans{$6$} & \ans{$4$} & \ans{$\langle 6,4\rangle$} \\
\hline
\end{tabular}

\vspace{0.1in}
\textbf{(b)} $\|\vec{p}(2)\| = \sqrt{6^{2}+4^{2}} = \sqrt{36+16} = \sqrt{52} = $ \ans{$2\sqrt{13}$}

\end{tcolorbox}

\vspace{0.1in}

%----------------------------------------------------------------------
% Tier A
%----------------------------------------------------------------------
\begin{tcolorbox}[colback=white, colframe=black!40,
    title=\textbf{Tier A --- Apply --- Key},
    fonttitle=\bfseries, arc=2mm, left=3mm, right=3mm, top=2mm, bottom=2mm]
\small

$f(t)=(t^{2}-4,\;3t)$.

\vspace{0.06in}
\textbf{(a)} $\vec{p}(t) = $ \ans{$\langle t^{2}-4,\;3t\rangle$}

\vspace{0.1in}
\textbf{(b)} $t^{2}-4=0$ $\Rightarrow$ $t=$ \ans{$\pm 2$}

Positions on $y$-axis: \ans{$(0,\,6)$ when $t=2$; $(0,\,-6)$ when $t=-2$}

\vspace{0.1in}
\textbf{(c)} $\vec{v}(1) = \langle 2,\;3\rangle$ \qquad
Speed $= \sqrt{4+9} = $ \ans{$\sqrt{13}$}

\begin{teachernote}
  Part (b): accept either order for the two crossings.  For part (c), students who
  used $t=1$ in $\vec{p}$ instead of applying the given rates have confused position
  and velocity --- redirect them to EK 4.9.A.2.
\end{teachernote}

\end{tcolorbox}

\vspace{0.1in}

%----------------------------------------------------------------------
% Tier E
%----------------------------------------------------------------------
\begin{tcolorbox}[colback=white, colframe=black!40,
    title=\textbf{Tier E --- Extend --- Key},
    fonttitle=\bfseries, arc=2mm, left=3mm, right=3mm, top=2mm, bottom=2mm]
\small

$\vec{p}(t)=\langle 2\cos t,\;2\sin t\rangle$.

\vspace{0.06in}
\textbf{(a)} \ans{$\|\vec{p}(t)\| = \sqrt{(2\cos t)^{2}+(2\sin t)^{2}}
= \sqrt{4\cos^{2}t+4\sin^{2}t}
= \sqrt{4(\cos^{2}t+\sin^{2}t)}
= \sqrt{4\cdot 1} = 2.$\\[3pt]
Therefore $\|\vec{p}(t)\|=2$ for all $t$ --- the particle moves on a circle of radius~2.}

\vspace{0.08in}
\textbf{(b)} $\vec{v}(0) = \langle -2\sin 0,\;2\cos 0\rangle = $ \ans{$\langle 0,\;2\rangle$}

\vspace{0.06in}
\textbf{(c)} At $t=0$ the particle is at \ans{$(2,\,0)$}.
Since $\vec{v}(0)=\langle 0,\,2\rangle$, the $x$-component is $0$ and the
$y$-component is $2>0$, so the particle is moving \ans{straight upward (in the
positive $y$-direction) at that instant}.

\end{tcolorbox}

\end{document}
